Weil's criterion expresses the Riemann hypothesis as the positivity of an explicit Hermitian form $\Q$ on compactly supported
tests on the multiplicative group. We fix Riemann's function $\Phi$, whose Fourier transform is the Riemann $\Xi$ function, in the normalisation of
Connes--Consani--Moscovici, and show that $f_0=\Phi/\|\Phi\|_2$, all its translates and all its convolutions lie in the radical of
$\Q$ on the natural enlarged test class. Dividing by $f_0$ gives a ground-state representation (in the sense of Frank--Seiringer) of Suzuki's
signed distributional kernel as a Dirichlet form, $\Q(f_0s)=\tfrac12\iint f_0(x)f_0(x')|s(x)-s(x')|^2\,d\tilde\nu(x-x')\,dx'$, with an explicit
measure: a density that changes sign exactly at $\log\pl$, $\pl$ the plastic number, and positive atoms $\Lambda(n)/\sqrt n$ at the prime-power
distances. The Riemann hypothesis becomes one explicit inequality with three terms: the prime atoms and the short-range positive density against
the long-range negative density. We prove that this inequality admits no local certificate: for any positive integrable function in the radical
of a functional with the cutoff property, every nonnegative minorant by weighted squares of finite difference stencils vanishes identically,
for stencils of any size and diameter; for the Weil form no sum-of-squares representation exists. The obstruction is the infinite-dimensional
radical, not the sign of the kernel, and it shows that the known nonlocal (Sonin-space) positivity cannot be localised. We also note that no
uniform margin on a dense family of tests can exist, record what certified finite-window computations (Connes--Consani--Moscovici matrices
with bottoms down to $10^{-220}$) do and do not show, and state the open problem in its reduced form. The core of the obstruction theorem and the
finite algebra of the window computations are formalised in Lean~4 against Mathlib; we disclose how automated reasoning tools were used.